\documentclass{amsart}
% For dealing with references we use the comment environment
\usepackage{verbatim}
\newenvironment{reference}{\comment}{\endcomment}
%\newenvironment{reference}{}{}
\newenvironment{slogan}{\comment}{\endcomment}
\newenvironment{history}{\comment}{\endcomment}

% For commutative diagrams we use Xy-pic
\usepackage[all]{xy}

% We use 2cell for 2-commutative diagrams.
\xyoption{2cell}
\UseAllTwocells

% We use multicol for the list of chapters between chapters
\usepackage{multicol}

% This is generally recommended for better output
\usepackage{lmodern}
\usepackage[T1]{fontenc}

% For cross-file-references

% Package for hypertext links:
\usepackage{hyperref}

% For any local file, say "hello.tex" you want to link to please

% Theorem environments.
%
\theoremstyle{plain}
\newtheorem{theorem}[subsection]{Theorem}
\newtheorem{proposition}[subsection]{Proposition}
\newtheorem{lemma}[subsection]{Lemma}

\theoremstyle{definition}
\newtheorem{definition}[subsection]{Definition}
\newtheorem{example}[subsection]{Example}
\newtheorem{exercise}[subsection]{Exercise}
\newtheorem{situation}[subsection]{Situation}

\theoremstyle{remark}
\newtheorem{remark}[subsection]{Remark}
\newtheorem{remarks}[subsection]{Remarks}

\numberwithin{equation}{subsection}

% Macros
%
\def\lim{\mathop{\mathrm{lim}}\nolimits}
\def\colim{\mathop{\mathrm{colim}}\nolimits}
\def\Spec{\mathop{\mathrm{Spec}}}
\def\Hom{\mathop{\mathrm{Hom}}\nolimits}
\def\Ext{\mathop{\mathrm{Ext}}\nolimits}
\def\SheafHom{\mathop{\mathcal{H}\!\mathit{om}}\nolimits}
\def\SheafExt{\mathop{\mathcal{E}\!\mathit{xt}}\nolimits}
\def\Sch{\mathit{Sch}}
\def\Mor{\mathop{\mathrm{Mor}}\nolimits}
\def\Ob{\mathop{\mathrm{Ob}}\nolimits}
\def\Sh{\mathop{\mathit{Sh}}\nolimits}
\def\NL{\mathop{N\!L}\nolimits}
\def\CH{\mathop{\mathrm{CH}}\nolimits}
\def\proetale{{pro\text{-}\acute{e}tale}}
\def\etale{{\acute{e}tale}}
\def\QCoh{\mathit{QCoh}}
\def\Ker{\mathop{\mathrm{Ker}}}
\def\Im{\mathop{\mathrm{Im}}}
\def\Coker{\mathop{\mathrm{Coker}}}
\def\Coim{\mathop{\mathrm{Coim}}}

% Boxtimes
%
\DeclareMathSymbol{\boxtimes}{\mathbin}{AMSa}{"02}

%
% Macros for moduli stacks/spaces
%
\def\QCohstack{\mathcal{QC}\!\mathit{oh}}
\def\Cohstack{\mathcal{C}\!\mathit{oh}}
\def\Spacesstack{\mathcal{S}\!\mathit{paces}}
\def\Quotfunctor{\mathrm{Quot}}
\def\Hilbfunctor{\mathrm{Hilb}}
\def\Curvesstack{\mathcal{C}\!\mathit{urves}}
\def\Polarizedstack{\mathcal{P}\!\mathit{olarized}}
\def\Complexesstack{\mathcal{C}\!\mathit{omplexes}}
% \Pic is the operator that assigns to X its picard group, usage \Pic(X)
% \Picardstack_{X/B} denotes the Picard stack of X over B
% \Picardfunctor_{X/B} denotes the Picard functor of X over B
\def\Pic{\mathop{\mathrm{Pic}}\nolimits}
\def\Picardstack{\mathcal{P}\!\mathit{ic}}
\def\Picardfunctor{\mathrm{Pic}}
\def\Deformationcategory{\mathcal{D}\!\mathit{ef}}

\setlength{\emergencystretch}{3em}
\begin{document}
\title[Stacks Project, Tag 0CRM]{1322 --- Countable cohomology over a cohomologically countable differential graded algebra is equivalent to a sequential homotopy colimit of compact objects}
\maketitle
\noindent Copyright (C) 2005--2025 Johan de Jong. Stacks Project authors. Source: \href{https://stacks.math.columbia.edu/tag/0CRM}{Tag 0CRM}.

\noindent Editorial difficulty note (not author text). Difficulty H2: The author constructs a countable quasi-isomorphic differential graded subalgebra by repeatedly killing all countable kernel classes, then a countable semi-free resolution and a sequential exhaustion by finite differential-closed cell submodules. The coupled algebra/model reduction and finite-cell closure construction give a substantive derived finiteness characterization beyond qualification homological algebra..

\noindent Editorial provenance note (not author text). History: excerpt from revision \texttt{a0444\allowbreak 6e57e\allowbreak c1fbc\allowbreak 252a8\allowbreak 71afc\allowbreak ec775\allowbreak 2fb28\allowbreak 07b14}, dga.tex, lines 7070--7199. Reference presentation was adapted; original-author context and core support blocks were retained or added. The mathematical wording is unchanged. Licensed under GNU FDL 1.2 or later; no invariant sections or cover texts. See ../licenses/COPYING. Context blocks reproduce author definitions and supporting results; they are not additional counted problems.

\begin{definition}
\label{definition-dga}
Let $R$ be a commutative ring. A {\it differential graded algebra over $R$}
is either
\begin{enumerate}
\item a chain complex $A_\bullet$ of $R$-modules endowed with
$R$-bilinear maps $A_n \times A_m \to A_{n + m}$,
$(a, b) \mapsto ab$ such that
$$
\text{d}_{n + m}(ab) = \text{d}_n(a)b + (-1)^n a\text{d}_m(b)
$$
and such that $\bigoplus A_n$ becomes an associative and unital
$R$-algebra, or
\item a cochain complex $A^\bullet$ of $R$-modules endowed with
$R$-bilinear maps $A^n \times A^m \to A^{n + m}$, $(a, b) \mapsto ab$
such that
$$
\text{d}^{n + m}(ab) = \text{d}^n(a)b + (-1)^n a\text{d}^m(b)
$$
and such that $\bigoplus A^n$ becomes an associative and unital $R$-algebra.
\end{enumerate}
\end{definition}

\begin{definition}
\label{definition-dgm}
Let $R$ be a ring.
Let $(A, \text{d})$ be a differential graded algebra over $R$.
A (right) {\it differential graded module} $M$ over $A$ is a right $A$-module
$M$ which has a grading $M = \bigoplus M^n$ and a differential $\text{d}$
such that $M^n A^m \subset M^{n + m}$, such that
$\text{d}(M^n) \subset M^{n + 1}$, and such that
$$
\text{d}(ma) = \text{d}(m)a + (-1)^n m\text{d}(a)
$$
for $a \in A$ and $m \in M^n$. A
{\it homomorphism of differential graded modules} $f : M \to N$
is an $A$-module map compatible with gradings and differentials.
The category of (right) differential graded $A$-modules is denoted
$\text{Mod}_{(A, \text{d})}$.
\end{definition}

\begin{definition}
\label{derived-definition-derived-colimit}
Let $\mathcal{D}$ be a triangulated category.
Let $(K_n, f_n)$ be a system of objects of $\mathcal{D}$.
We say an object $K$ is a {\it derived colimit}, or a
{\it homotopy colimit} of the system $(K_n)$ if
the direct sum $\bigoplus K_n$ exists and there is a distinguished triangle
$$
\bigoplus K_n \to \bigoplus K_n \to K \to \bigoplus K_n[1]
$$
where the map $\bigoplus K_n \to \bigoplus K_n$ is given
by $1 - f_n$ in degree $n$. If this is the
case, then we sometimes indicate this by the notation
$K = \text{hocolim} K_n$.
\end{definition}

\begin{lemma}
\label{lemma-resolve}
Let $(A, \text{d})$ be a differential graded algebra.
Let $M$ be a differential graded $A$-module. There exists a homomorphism
$P \to M$ of differential graded $A$-modules such that
\begin{enumerate}
\item $P \to M$ is a quasi-isomorphism, and
\item $P$ has property (P).
\end{enumerate}
\end{lemma}

\begin{proof}
Set $M = M_0$. We inductively choose short exact sequences
$$
0 \to M_{i + 1} \to P_i \to M_i \to 0
$$
where the maps $P_i \to M_i$ are chosen as in Lemma \ref{lemma-good-quotient}.
This gives a ``resolution''
$$
\ldots \to P_2 \xrightarrow{f_2} P_1 \xrightarrow{f_1} P_0 \to M \to 0
$$
Then we set
$$
P = \bigoplus\nolimits_{i \geq 0} P_i
$$
as an $A$-module with grading given by
$P^n = \bigoplus_{a + b = n} P_{-a}^b$ and
differential (as in the construction of the total complex associated
to a double complex) by
$$
\text{d}_P(x) = f_{-a}(x) + (-1)^a \text{d}_{P_{-a}}(x)
$$
for $x \in P_{-a}^b$. With these conventions $P$ is indeed a differential
graded $A$-module. Recalling that each $P_i$ has a two step filtration
$0 \to P_i' \to P_i \to P_i'' \to 0$ we set
$$
F_{2i}P = \bigoplus\nolimits_{i \geq j \geq 0} P_j
\subset
\bigoplus\nolimits_{i \geq 0} P_i = P
$$
and we add $P'_{i + 1}$ to $F_{2i}P$ to get $F_{2i + 1}$.
These are differential graded submodules and the successive quotients
are direct sums of shifts of $A$. By
Lemma \href{https://stacks.math.columbia.edu/tag/09K0}{lemma target graded projective (Tag 09K0)} we see that
the inclusions $F_iP \to F_{i + 1}P$ are admissible monomorphisms.
Finally, we have to show that the map $P \to M$ (given by the
augmentation $P_0 \to M$) is a quasi-isomorphism. This follows from
Homology, Lemma \href{https://stacks.math.columbia.edu/tag/09IZ}{lemma good resolution gives qis (Tag 09IZ)}.
\end{proof}


\begin{lemma}
\label{lemma-good-quotient}
Let $(A, \text{d})$ be a differential graded algebra.
Let $M$ be a differential graded $A$-module. There exists a homomorphism
$P \to M$ of differential graded $A$-modules with the following
properties
\begin{enumerate}
\item $P \to M$ is surjective,
\item $\Ker(\text{d}_P) \to \Ker(\text{d}_M)$ is surjective, and
\item $P$ sits in an admissible short exact sequence
$0 \to P' \to P \to P'' \to 0$ where $P'$, $P''$ are direct sums
of shifts of $A$.
\end{enumerate}
\end{lemma}

\begin{proof}
Let $P_k$ be the free $A$-module with generators $x, y$ in degrees
$k$ and $k + 1$. Define the structure of a differential graded
$A$-module on $P_k$ by setting $\text{d}(x) = y$ and $\text{d}(y) = 0$.
For every element $m \in M^k$ there is a homomorphism
$P_k \to M$ sending $x$ to $m$ and $y$ to $\text{d}(m)$.
Thus we see that there is a surjection from a direct sum
of copies of $P_k$ to $M$. This clearly produces $P \to M$
having properties (1) and (3). To obtain property (2) note
that if $m \in \Ker(\text{d}_M)$ has degree $k$, then there is a map
$A[k] \to M$ mapping $1$ to $m$. Hence we can achieve (2) by adding
a direct sum of copies of shifts of $A$.
\end{proof}


\begin{proposition}
\label{proposition-compact}
Let $(A, \text{d})$ be a differential graded algebra. Let $E$ be an
object of $D(A, \text{d})$. Then the following are equivalent
\begin{enumerate}
\item $E$ is a compact object,
\item $E$ is a direct summand of an object of $D(A, \text{d})$
which is represented by a differential graded module $P$ which
has a finite filtration $F_\bullet$ by differential graded submodules
such that $F_iP/F_{i - 1}P$ are finite direct sums of shifts of $A$.
\end{enumerate}
\end{proposition}

\begin{proof}
Assume $E$ is compact. By Lemma \ref{lemma-resolve} we may assume that $E$
is represented by a differential graded $A$-module $P$ with property (P).
Consider the distinguished triangle
$$
\bigoplus F_iP \to \bigoplus F_iP \to P
\xrightarrow{\delta} \bigoplus F_iP[1]
$$
coming from the admissible short exact sequence of
Lemma \ref{lemma-property-P-sequence}. Since $E$ is compact we have
$\delta = \sum_{i = 1, \ldots, n} \delta_i$ for some
$\delta_i : P \to F_iP[1]$. Since the composition of $\delta$
with the map $\bigoplus F_iP[1] \to \bigoplus F_iP[1]$ is zero
(Derived Categories, Lemma \href{https://stacks.math.columbia.edu/tag/0146}{lemma composition zero (Tag 0146)})
it follows that $\delta = 0$ (follows as $\bigoplus F_iP \to \bigoplus F_iP$
maps the summand $F_iP$ via the difference of $\text{id}$ and the inclusion
map into $F_{i - 1}P$).
Thus we see that the identity on $E$ factors through
$\bigoplus F_iP$ in $D(A, \text{d})$ (by
Derived Categories, Lemma \href{https://stacks.math.columbia.edu/tag/05QT}{lemma split (Tag 05QT)}).
Next, we use that $P$ is compact again to see that the map
$E \to \bigoplus F_iP$ factors through $\bigoplus_{i = 1, \ldots, n} F_iP$
for some $n$. In other words, the identity on $E$ factors through
$\bigoplus_{i = 1, \ldots, n} F_iP$. By
Lemma \ref{lemma-factor-through-nicer}
we see that the identity of $E$ factors as $E \to P \to E$
where $P$ is as in part (2) of the statement of the lemma.
In other words, we have proven that (1) implies (2).

\medskip\noindent
Assume (2). By
Derived Categories, Lemma \href{https://stacks.math.columbia.edu/tag/09QH}{lemma compact objects subcategory (Tag 09QH)}
it suffices to show that $P$ gives a compact object. Observe that
$P$ has property (P), hence we have
$$
\Hom_{D(A, \text{d})}(P, M) = \Hom_{K(A, \text{d})}(P, M)
$$
for any differential graded module $M$ by Lemma \ref{lemma-hom-derived}.
As direct sums in $D(A, \text{d})$ are given by direct sums of
graded modules (Lemma \href{https://stacks.math.columbia.edu/tag/09QI}{lemma derived products (Tag 09QI)}) we reduce
to showing that $\Hom_{K(A, \text{d})}(P, M)$ commutes with direct
sums. Using that $K(A, \text{d})$ is a triangulated category,
that $\Hom$ is a cohomological functor in the first
variable, and the filtration on $P$, we reduce to the case that
$P$ is a finite direct sum of shifts of $A$. Thus we reduce to
the case $P = A[k]$ which is clear.
\end{proof}


\begin{lemma}
\label{lemma-factor-through-nicer}
Let $(A, \text{d})$ be a differential graded algebra. Let $E$ be a compact
object of $D(A, \text{d})$. Let $P$ be a differential graded $A$-module
which has a finite filtration
$$
0 = F_{-1}P \subset F_0P \subset F_1P \subset \ldots \subset F_nP = P
$$
by differential graded submodules such that
$$
F_{i + 1}P/F_iP \cong \bigoplus\nolimits_{j \in J_i} A[k_{i, j}]
$$
as differential graded $A$-modules for some sets $J_i$ and integers $k_{i, j}$.
Let $E \to P$ be a morphism of $D(A, \text{d})$.
Then there exists a differential graded submodule $P' \subset P$ such that
$F_{i + 1}P \cap P'/(F_iP \cap P')$ is equal to
$\bigoplus_{j \in J'_i} A[k_{i, j}]$ for some finite subsets
$J'_i \subset J_i$ and such that $E \to P$ factors through $P'$.
\end{lemma}

\begin{proof}
We will prove by induction on $-1 \leq m \leq n$ that there exists
a differential graded submodule $P' \subset P$  such that
\begin{enumerate}
\item $F_mP \subset P'$,
\item for $i \geq m$ the quotient $F_{i + 1}P \cap P'/(F_iP \cap P')$ is
isomorphic to $\bigoplus_{j \in J'_i} A[k_{i, j}]$ for some finite subsets
$J'_i \subset J_i$, and
\item $E \to P$ factors through $P'$.
\end{enumerate}
The base case is $m = n$ where we can take $P' = P$.

\medskip\noindent
Induction step. Assume $P'$ works for $m$.
For $i \geq m$ and $j \in J'_i$ let $x_{i, j} \in F_{i + 1}P \cap P'$
be a homogeneous element of degree $k_{i, j}$ whose image in
$F_{i + 1}P \cap P'/(F_iP \cap P')$ is the generator in
the summand corresponding to $j \in J_i$. The
$x_{i, j}$ generate $P'/F_mP$ as an $A$-module. Write
$$
\text{d}(x_{i, j}) = \sum x_{i', j'} a_{i, j}^{i', j'} + y_{i, j}
$$
with $y_{i, j} \in F_mP$ and $a_{i, j}^{i', j'} \in A$.
There exists a finite subset
$J'_{m - 1} \subset J_{m - 1}$ such that each $y_{i, j}$ maps to
an element of the submodule $\bigoplus_{j \in J'_{m - 1}} A[k_{m - 1, j}]$
of $F_mP/F_{m - 1}P$. Let $P'' \subset F_mP$ be the inverse
image of $\bigoplus_{j \in J'_{m - 1}} A[k_{m - 1, j}]$ under
the map $F_mP \to F_mP/F_{m - 1}P$. Then we see that the $A$-submodule
$$
P'' + \sum x_{i, j}A
$$
is a differential graded submodule of the type we are looking for. Moreover
$$
P'/(P'' + \sum x_{i, j}A) =
\bigoplus\nolimits_{j \in J_{m - 1} \setminus J'_{m - 1}} A[k_{m - 1, j}]
$$
Since $E$ is compact, the composition of the given map $E \to P'$
with the quotient map, factors through a finite direct subsum of
the module displayed above. Hence after enlarging $J'_{m - 1}$
we may assume $E \to P'$ factors through
$P'' + \sum x_{i, j}A$ as desired.
\end{proof}


\begin{lemma}
\label{lemma-property-P-sequence}
Let $(A, \text{d})$ be a differential graded algebra.
Let $P$ be a differential graded $A$-module. If $F_\bullet$
is a filtration as in property (P), then we obtain an
admissible short exact sequence
$$
0 \to
\bigoplus\nolimits F_iP \to
\bigoplus\nolimits F_iP \to P \to 0
$$
of differential graded $A$-modules.
\end{lemma}

\begin{proof}
The second map is the direct sum of the inclusion maps.
The first map on the summand $F_iP$ of the source is the sum
of the identity $F_iP \to F_iP$ and the negative of the inclusion
map $F_iP \to F_{i + 1}P$. Choose homomorphisms $s_i : F_{i + 1}P \to F_iP$
of graded $A$-modules which are left inverse to the inclusion maps.
Composing gives maps $s_{j, i} : F_jP \to F_iP$ for all $j > i$.
Then a left inverse of the first arrow maps $x \in F_jP$ to
$(s_{j, 0}(x), s_{j, 1}(x), \ldots, s_{j, j - 1}(x), 0, \ldots)$
in $\bigoplus F_iP$.
\end{proof}


\begin{lemma}
\label{lemma-property-P-K-projective}
Let $(A, \text{d})$ be a differential graded algebra.
Let $P$ be a differential graded $A$-module with property (P).
Then
$$
\Hom_{K(\text{Mod}_{(A, \text{d})})}(P, N) = 0
$$
for all acyclic differential graded $A$-modules $N$.
\end{lemma}

\begin{proof}
We will use that $K(\text{Mod}_{(A, \text{d})})$ is a triangulated
category (Proposition \href{https://stacks.math.columbia.edu/tag/09KJ}{proposition homotopy category triangulated (Tag 09KJ)}).
Let $F_\bullet$ be a filtration on $P$ as in property (P).
The short exact sequence of Lemma \ref{lemma-property-P-sequence}
produces a distinguished triangle. Hence by
Derived Categories, Lemma \href{https://stacks.math.columbia.edu/tag/0149}{lemma representable homological (Tag 0149)}
it suffices to show that
$$
\Hom_{K(\text{Mod}_{(A, \text{d})})}(F_iP, N) = 0
$$
for all acyclic differential graded $A$-modules $N$ and all $i$.
Each of the differential graded modules $F_iP$ has a finite filtration
by admissible monomorphisms, whose graded pieces are direct sums
of shifts $A[k]$. Thus it suffices to prove that
$$
\Hom_{K(\text{Mod}_{(A, \text{d})})}(A[k], N) = 0
$$
for all acyclic differential graded $A$-modules $N$ and all $k$.
This follows from Lemma \href{https://stacks.math.columbia.edu/tag/09K1}{lemma hom from shift free (Tag 09K1)}.
\end{proof}


\begin{lemma}
\label{lemma-hom-derived}
Let $(A, \text{d})$ be a differential graded algebra.
Let $M$ and $N$ be differential graded $A$-modules.
\begin{enumerate}
\item Let $P \to M$ be a P-resolution as in
Lemma \ref{lemma-resolve}. Then
$$
\Hom_{D(A, \text{d})}(M, N) =
\Hom_{K(\text{Mod}_{(A, \text{d})})}(P, N)
$$
\item Let $N \to I$ be an I-resolution as in
Lemma \href{https://stacks.math.columbia.edu/tag/09KU}{lemma right resolution (Tag 09KU)}. Then
$$
\Hom_{D(A, \text{d})}(M, N) =
\Hom_{K(\text{Mod}_{(A, \text{d})})}(M, I)
$$
\end{enumerate}
\end{lemma}

\begin{proof}
Let $P \to M$ be as in (1). Since $P \to M$ is a quasi-isomorphism we see that
$$
\Hom_{D(A, \text{d})}(P, N) = \Hom_{D(A, \text{d})}(M, N)
$$
by definition of the derived category. A morphism
$f : P \to N$ in $D(A, \text{d})$ is equal to
$s^{-1}f'$ where $f' : P \to N'$ is a morphism and
$s : N \to N'$ is a quasi-isomorphism. Choose a distinguished triangle
$$
N \to N' \to Q \to N[1]
$$
As $s$ is a quasi-isomorphism, we see that $Q$ is acyclic. Thus
$\Hom_{K(\text{Mod}_{(A, \text{d})})}(P, Q[k]) = 0$ for all $k$ by
Lemma \ref{lemma-property-P-K-projective}. Since
$\Hom_{K(\text{Mod}_{(A, \text{d})})}(P, -)$
is cohomological, we conclude that we can lift $f' : P \to N'$
uniquely to a morphism $f : P \to N$. This finishes the proof.

\medskip\noindent
The proof of (2) is dual to that of (1) using
Lemma \href{https://stacks.math.columbia.edu/tag/09KS}{lemma property I K injective (Tag 09KS)} in stead of
Lemma \ref{lemma-property-P-K-projective}.
\end{proof}


\begin{lemma}
\label{lemma-homotopy-colimit}
Let $(A, \text{d})$ be a differential graded algebra. Let
$M_n$ be a system of differential graded modules. Then the derived
colimit $\text{hocolim} M_n$ in $D(A, \text{d})$ is represented
by the differential graded module $\colim M_n$.
\end{lemma}

\begin{proof}
Set $M = \colim M_n$. We have an exact sequence of differential graded modules
$$
0 \to \bigoplus M_n \to \bigoplus M_n \to M \to 0
$$
by Derived Categories, Lemma \href{https://stacks.math.columbia.edu/tag/093W}{lemma compute colimit (Tag 093W)}
(applied the underlying complexes of abelian groups).
The direct sums are direct sums in $D(\mathcal{A})$ by
Lemma \href{https://stacks.math.columbia.edu/tag/09QI}{lemma derived products (Tag 09QI)}.
Thus the result follows from the definition
of derived colimits in Derived Categories,
Definition \ref{derived-definition-derived-colimit}
and the fact that a short exact sequence of complexes
gives a distinguished triangle
(Lemma \href{https://stacks.math.columbia.edu/tag/09L0}{lemma derived canonical delta functor (Tag 09L0)}).
\end{proof}


\section{P-resolutions}
\label{section-P-resolutions}

\noindent
This section is the analogue of
Derived Categories, Section \href{https://stacks.math.columbia.edu/tag/06XW}{section unbounded (Tag 06XW)}.

\medskip\noindent
Let $(A, \text{d})$ be a differential graded algebra.
Let $P$ be a differential graded $A$-module. We say $P$
{\it has property (P)} if it there exists a filtration
$$
0 = F_{-1}P \subset F_0P \subset F_1P \subset \ldots \subset P
$$
by differential graded submodules such that
\begin{enumerate}
\item $P = \bigcup F_pP$,
\item the inclusions $F_iP \to F_{i + 1}P$ are admissible
monomorphisms,
\item the quotients $F_{i + 1}P/F_iP$ are isomorphic as differential
graded $A$-modules to a direct sum of $A[k]$.
\end{enumerate}
In fact, condition (2) is a consequence of condition (3), see
Lemma \href{https://stacks.math.columbia.edu/tag/09K0}{lemma target graded projective (Tag 09K0)}. Moreover, the reader
can verify that as a graded $A$-module $P$ will be isomorphic to a
direct sum of shifts of $A$.


\begin{lemma}
\label{lemma-countable}
Let $(A, \text{d})$ be a differential graded algebra with
$H^i(A)$ countable for each $i$. Let $M$ be an object of $D(A, \text{d})$.
Then the following are equivalent
\begin{enumerate}
\item $M = \text{hocolim} E_n$ with $E_n$ compact in $D(A, \text{d})$, and
\item $H^i(M)$ is countable for each $i$.
\end{enumerate}
\end{lemma}

\begin{proof}
Assume (1) holds. Then we have $H^i(M) = \colim H^i(E_n)$ by
Derived Categories, Lemma \href{https://stacks.math.columbia.edu/tag/0CRK}{lemma cohomology of hocolim (Tag 0CRK)}.
Thus it suffices to prove that $H^i(E_n)$ is countable for each $n$.
By Proposition \ref{proposition-compact} we see that $E_n$
is isomorphic in $D(A, \text{d})$ to a direct summand of a
differential graded module $P$ which has a finite filtration
$F_\bullet$ by differential graded submodules
such that $F_jP/F_{j - 1}P$ are finite direct sums of shifts of $A$.
By assumption the groups $H^i(F_jP/F_{j - 1}P)$ are countable.
Arguing by induction on the length of the filtration and using
the long exact cohomology sequence we conclude that (2) is true.
The interesting implication is the other one.

\medskip\noindent
We claim there is a countable differential graded
subalgebra $A' \subset A$ such that the inclusion map
$A' \to A$ defines an isomorphism on cohomology.
To construct $A'$ we choose countable differential graded
subalgebras
$$
A_1 \subset A_2 \subset A_3 \subset \ldots
$$
such that (a) $H^i(A_1) \to H^i(A)$ is surjective, and (b)
for $n > 1$ the kernel of the map $H^i(A_{n - 1}) \to H^i(A_n)$
is the same as the kernel of the map $H^i(A_{n - 1}) \to H^i(A)$.
To construct $A_1$ take any countable collection of cochains
$S \subset A$ generating the cohomology of $A$ (as a ring or as
a graded abelian group) and let
$A_1$ be the differential graded subalgebra of $A$ generated by $S$.
To construct $A_n$ given $A_{n - 1}$ for each cochain $a \in A_{n - 1}^i$
which maps to zero in $H^i(A)$ choose $s_a \in A^{i - 1}$
with $\text{d}(s_a) = a$ and let $A_n$ be the differential graded
subalgebra of $A$ generated by $A_{n - 1}$ and the elements $s_a$.
Finally, take $A' = \bigcup A_n$.

\medskip\noindent
By Lemma \href{https://stacks.math.columbia.edu/tag/09S6}{lemma qis equivalence (Tag 09S6)}
the restriction map $D(A, \text{d}) \to D(A', \text{d})$,
$M \mapsto M_{A'}$ is an equivalence. Since the cohomology
groups of $M$ and $M_{A'}$ are the same, we see that it
suffices to prove the implication (2) $\Rightarrow$ (1)
for $(A', \text{d})$.

\medskip\noindent
Assume $A$ is countable. By the exact same type of argument as
given above we see that for $M$ in $D(A, \text{d})$
the following are equivalent: $H^i(M)$ is countable for each $i$
and $M$ can be represented by a countable differential graded module.
Hence in order to prove the implication (2) $\Rightarrow$ (1)
we reduce to the situation described in the next paragraph.

\medskip\noindent
Assume $A$ is countable and that $M$ is a countable differential graded
module over $A$. We claim there exists a homomorphism
$P \to M$ of differential graded $A$-modules such that
\begin{enumerate}
\item $P \to M$ is a quasi-isomorphism,
\item $P$ has property (P), and
\item $P$ is countable.
\end{enumerate}
Looking at the proof of the construction of P-resolutions in
Lemma \ref{lemma-resolve} we see that it suffices to show that
we can prove Lemma \ref{lemma-good-quotient}
in the setting of countable differential graded modules.
This is immediate from the proof.

\medskip\noindent
Assume that $A$ is countable and that $M$ is a countable
differential graded module with property (P). Choose a filtration
$$
0 = F_{-1}P \subset F_0P \subset F_1P \subset \ldots \subset P
$$
by differential graded submodules such that we have
\begin{enumerate}
\item $P = \bigcup F_pP$,
\item $F_iP \to F_{i + 1}P$ is an admissible monomorphism,
\item isomorphisms of differential graded modules
$F_iP/F_{i - 1}P \to \bigoplus_{j \in J_i} A[k_j]$
for some sets $J_i$ and integers $k_j$.
\end{enumerate}
Of course $J_i$ is countable for each $i$. For each $i$ and
$j \in J_i$ choose $x_{i, j} \in F_iP$ of degree $k_j$ whose
image in $F_iP/F_{i - 1}P$ generates the summand corresponding
to $j$.

\medskip\noindent
Claim: Given $n$ and finite subsets $S_i \subset J_i$, $i = 1, \ldots, n$
there exist finite subsets $S_i \subset T_i \subset J_i$, $i = 1, \ldots, n$
such that $P' = \bigoplus_{i \leq n} \bigoplus_{j \in T_i} Ax_{i, j}$
is a differential graded submodule of $P$. This was shown in the
proof of Lemma \ref{lemma-factor-through-nicer} but it is also
easily shown directly: the elements $x_{i, j}$ freely generate
$P$ as a right $A$-module. The structure of $P$ shows that
$$
\text{d}(x_{i, j}) = \sum\nolimits_{i' < i} x_{i', j'}a_{i', j'}
$$
where of course the sum is finite.
Thus given $S_0, \ldots, S_n$ we can first choose
$S_0 \subset S'_0, \ldots, S_{n - 1} \subset S'_{n - 1}$ with
$\text{d}(x_{n, j}) \in \bigoplus_{i' < n, j' \in S'_{i'}} x_{i', j'}A$
for all $j \in S_n$. Then by induction on $n$ we can choose
$S'_0 \subset T_0, \ldots, S'_{n - 1} \subset T_{n - 1}$
to make sure that $\bigoplus_{i' < n, j' \in T_{i'}} x_{i', j'}A$
is a differential graded $A$-submodule. Setting $T_n = S_n$ we find that
$P' = \bigoplus_{i \leq n, j \in T_i} x_{i, j}A$ is as desired.

\medskip\noindent
From the claim it is clear that $P = \bigcup P'_n$
is a countable rising union of $P'_n$ as above.
By construction each $P'_n$ is a differential graded module with
property (P) such that the filtration is finite and the successive
quotients are finite direct sums of shifts of $A$. Hence $P'_n$
defines a compact object of $D(A, \text{d})$, see for example
Proposition \ref{proposition-compact}. Since
$P = \text{hocolim} P'_n$ in $D(A, \text{d})$
by Lemma \ref{lemma-homotopy-colimit}
the proof of the implication (2) $\Rightarrow$ (1) is complete.
\end{proof}
\end{document}
